\section{Marco teórico}
~\label{sec:marco-teorico}

\subsection{Modelo OSI y arquitectura TCP/IP}

El diseño y operación de sistemas de monitoreo de red se sostienen 
sobre dos modelos de referencia que describen la comunicación entre 
dispositivos de forma estratificada: el modelo OSI de siete capas 
propuesto por la Organización Internacional de Normalización 
~\cite{iso7498}, y la arquitectura TCP/IP de cuatro capas que refleja 
la implementación práctica de la suite de protocolos de Internet 
~\cite{forouzan2006transmision}. Ambos modelos abstraen la complejidad 
de las comunicaciones mediante la separación de responsabilidades: 
cada capa ofrece un servicio específico a la capa inmediatamente 
superior y se apoya en los servicios de la capa inmediatamente 
inferior, lo cual permite que los componentes de un sistema 
evolucionen de forma independiente siempre que respeten las 
interfaces entre capas~\cite{kurose2010redes}.

Las capas inferiores determinan los mecanismos sobre los cuales opera 
GOATGuard. La capa de enlace de datos (Capa 2) define la estructura 
de las tramas Ethernet~\cite{ieee8023} y las direcciones MAC, que 
permiten la identificación de dispositivos a nivel físico dentro del 
segmento local; esta capa soporta el descubrimiento de activos 
mediante el protocolo ARP~\cite{rfc826}. La capa de red (Capa 3) 
define las direcciones IP y los mecanismos de enrutamiento; sobre 
esta capa opera el protocolo ICMP~\cite{rfc792} empleado en el 
sondeo de salud del enlace hacia el proveedor de servicios. La capa 
de transporte (Capa 4) proporciona los dos protocolos que estructuran 
el transporte dual del sistema: TCP~\cite{rfc793} para la entrega 
confiable y ordenada del tráfico capturado, y UDP~\cite{rfc768} para 
la transmisión de métricas y señales de disponibilidad. Finalmente, 
la capa de aplicación (Capa 7) aloja los protocolos DNS 
~\cite{rfc1035}, HTTP~\cite{rfc9110} y los intercambios TLS 
~\cite{rfc8446} observados durante la captura, cuya inspección 
alimenta el análisis de flujos.

Aunque los dos modelos difieren en el número de capas, las funciones 
relevantes para este trabajo coinciden en ambos. Esto permite 
expresar las decisiones de diseño en términos de capas sin ambigüedad 
y asociar cada componente del sistema a una responsabilidad 
específica dentro del stack de comunicaciones~\cite{kurose2010redes}.

\subsection{Estructura de un paquete de red
}

Un paquete observado en una interfaz Ethernet sigue un esquema de 
encapsulación anidada en el que cada capa antepone sus propios 
encabezados al contenido recibido de la capa superior 
\cite{forouzan2006transmision}. En la base se encuentra el encabezado 
Ethernet de catorce bytes~\cite{ieee8023}, que incluye las 
direcciones MAC de origen y destino y el campo EtherType que indica 
el protocolo de la capa siguiente (Fig.~\ref{fig:paquete-red}). A 
continuación aparece el encabezado IP de veinte bytes en su versión 
mínima, que contiene las direcciones IP de origen y destino, el 
identificador de protocolo de transporte, el TTL y el checksum de 
cabecera. Por último se ubica el encabezado de transporte: veinte 
bytes para TCP, con números de secuencia y acuse, tamaño de ventana 
y banderas de control, u ocho bytes para UDP, con puertos y longitud 
del datagrama~\cite{stevens2004unix}.

Este conjunto de encabezados, aproximadamente cincuenta y cuatro 
bytes, concentra la totalidad de la información necesaria para 
reconstruir flujos y extraer métricas de comportamiento sin recurrir 
al contenido del payload. Las direcciones MAC identifican emisor y 
receptor dentro del segmento local; las direcciones IP, junto con 
los puertos y el protocolo de transporte, conforman la tupla de 
cinco elementos que define cada flujo de comunicación 
~\cite{kurose2010redes}; las banderas TCP permiten determinar el 
ciclo de vida de una conexión y detectar retransmisiones 
~\cite{rfc793}. El payload de aplicación, por el contrario, puede 
contener información sensible como credenciales, contenido de 
comunicaciones o datos personales, y su captura integral 
comprometería la confidencialidad de los usuarios monitoreados.
Un segundo aspecto relevante es el orden de los bytes en la 
transmisión. Los encabezados de los protocolos de red se representan 
en network byte order, que corresponde a la convención big-endian: 
el byte más significativo se transmite primero~\cite{rfc793}. Los 
sistemas operativos sobre arquitecturas x86 almacenan los datos en 
little-endian, por lo que la serialización y deserialización de 
campos numéricos requiere conversiones explícitas para garantizar la 
interoperabilidad entre emisores y receptores independientemente de 
la arquitectura subyacente~\cite{stevens2004unix}.


\begin{figure}[H]
\centering
\includegraphics[width=\columnwidth]{paquete_red.png}
\caption{Estructura simplificada de un paquete de red.}
~\label{fig:paquete-red}
\end{figure}

Esta estructura justifica técnicamente la estrategia de sanitización 
adoptada por los sensores de captura: un recorte de noventa y seis 
bytes preserva los encabezados de las tres capas inferiores con 
margen para las opciones TCP, mientras que un recorte extendido a 
trescientos bytes para tráfico dirigido a los puertos asociados a 
DNS y HTTPS permite incluir el dominio consultado y la indicación de 
nombre de servidor del handshake TLS~\cite{rfc8446}, elementos con 
valor diagnóstico que no comprometen el contenido de las 
comunicaciones.

\subsection{Monitoreo de redes: captura centralizada vs distribuida}

Los enfoques tradicionales de monitoreo en redes de área local 
operan desde un único punto de observación, típicamente un servidor 
que recibe una copia del tráfico mediante un puerto mirror en el 
switch o una interfaz configurada para escuchar el segmento 
~\cite{kurose2010redes}. Sobre este modelo pasivo se articulan el 
monitoreo activo, que genera tráfico de prueba para verificar 
conectividad~\cite{forouzan2006transmision}, y herramientas de uso 
extendido como Wireshark para análisis de paquetes, Nmap para 
descubrimiento de hosts~\cite{lyon2009nmap}, Snort para detección de 
intrusos~\cite{roesch1999snort}, y plataformas empresariales como 
Zabbix o PRTG para consolidación de métricas. La limitación 
estructural de este enfoque es la pérdida de granularidad: la 
captura agregada impide atribuir un patrón anómalo al endpoint que 
lo origina, y su dependencia de consolas fijas, sumada al 
licenciamiento orientado a grandes infraestructuras, lo hace inviable 
en redes de pequeña escala.

La alternativa distribuida despliega procesos ligeros en cada 
endpoint que capturan el tráfico localmente y lo reportan a un 
servidor central~\cite{paxson1999zeek}. Esta topología preserva la 
granularidad por dispositivo y habilita la correlación directa entre 
una anomalía y su equipo responsable, a costa de instalar un 
componente en cada host y diseñar un canal de transporte que no 
sature la red ni degrade el rendimiento del endpoint. Estas 
tensiones de diseño orientan las decisiones arquitectónicas 
descritas en la Sección~\ref{sec:desarrollo}.

El modelo distribuido cuenta con un antecedente formal en la 
arquitectura IPFIX (IP Flow Information Export) especificada por la 
IETF en la RFC 5470~\cite{rfc5470}. Este documento establece un 
modelo de referencia en el que uno o varios \textit{observation 
points} dentro de un \textit{IPFIX device} alimentan un conjunto de 
\textit{metering processes} que a su vez entregan sus datos a un 
\textit{exporting process}, el cual los transmite a un 
\textit{collector} encargado de la recepción, procesamiento y 
almacenamiento (Fig.~\ref{fig:ipfix}). La correspondencia con 
GOATGuard es directa: los sensores de captura cumplen el rol de 
observation point y metering process sobre cada endpoint, mientras 
que el backend actúa como collector que consolida, procesa y 
persiste los datos reportados. Esta correspondencia no implica 
adoptar el protocolo IPFIX como tal sino reconocer que la 
descomposición funcional del sistema se apoya en un modelo de 
referencia ya estandarizado.

\begin{figure}[H]
\centering
\includegraphics[width=\columnwidth]{ipfix_rfc5470.png}
\caption{Modelo de referencia IPFIX definido en la RFC 5470~\cite{rfc5470}. La arquitectura de GOATGuard mapea sus sensores al rol de observation point y metering process, y el backend al rol de collector.}
~\label{fig:ipfix}
\end{figure}

\subsection{Descubrimiento y clasificación de activos}

El descubrimiento de dispositivos parte del escaneo activo: 
solicitudes ARP~\cite{rfc826} para resolver direcciones IP y MAC 
dentro del segmento de enlace, sondeos ICMP~\cite{rfc792} para 
confirmar alcanzabilidad, y pruebas a puertos TCP/UDP específicos 
para revelar servicios expuestos; herramientas como Nmap combinan 
estas técnicas en un inventario inicial~\cite{lyon2009nmap}. Dos 
limitaciones acotan este enfoque. Primero, la información recolectada 
se reduce a presencia y datos técnicos básicos, sin distinguir entre 
una estación de trabajo, una impresora, un dispositivo IoT o un 
equipo personal no autorizado. Segundo, los escaneos se ejecutan de 
forma periódica y no continua, dejando ventanas durante las cuales 
un dispositivo puede conectarse sin ser detectado hasta el siguiente 
ciclo.

La técnica de network fingerprinting responde parcialmente a la 
primera limitación combinando señales de múltiples capas del stack 
~\cite{lippmann2003fingerprint}: el prefijo OUI de la MAC identifica 
al fabricante, los puertos abiertos y banners de servicios revelan 
el sistema operativo y las aplicaciones, y el patrón de 
comunicación, como la frecuencia de consultas DNS o los destinos 
recurrentes, permite inferir el rol operativo del dispositivo. La 
combinación de estos indicadores habilita la transición desde un 
listado plano hacia un inventario clasificado. El sistema propuesto 
integra el descubrimiento ARP ejecutado por los sensores con el 
registro explícito de los endpoints instrumentados, manteniendo 
trazabilidad tanto de los dispositivos bajo control directo como de 
los observados pasivamente en el segmento.

\subsection{Detección de anomalías por métodos estadísticos adaptativos}

La detección de anomalías en series temporales consiste en 
identificar observaciones que se desvían significativamente del 
comportamiento esperado dado el historial previo de la misma 
variable~\cite{chandola2009anomaly}. En el contexto del monitoreo 
de infraestructura, cada métrica asociada a un dispositivo produce 
una secuencia de valores en el tiempo cuya distribución no es 
conocida a priori y varía entre dispositivos: una estación de 
desarrollo puede presentar un consumo de CPU sostenido del sesenta 
por ciento durante compilaciones, mientras que un servidor de 
archivos puede operar establemente por debajo del diez por ciento. 
Un método de detección útil debe construir un perfil individual 
para cada dispositivo y métrica, adaptarse a cambios legítimos en 
el comportamiento, y operar en línea sin requerir un conjunto de 
entrenamiento previo~\cite{chandola2009anomaly}.

Los umbrales fijos constituyen el enfoque más simple pero resultan 
inadecuados para esta tarea: un límite absoluto definido globalmente 
ignora el perfil particular de cada dispositivo y produce falsos 
positivos en los que operan legítimamente cerca del umbral, al 
tiempo que pierde anomalías relativas en dispositivos con líneas 
base bajas~\cite{montgomery2013spc}. El cálculo de Z-Score sobre la 
media aritmética acumulada corrige parcialmente este problema al 
normalizar por la desviación típica observada, pero la media 
acumulada pondera por igual todas las observaciones históricas y 
tarda horas en converger cuando el comportamiento del dispositivo 
cambia de forma sostenida, lo que genera falsas alarmas durante el 
periodo de adaptación.

La media móvil ponderada exponencialmente (Exponentially Weighted 
Moving Average, EWMA) propuesta originalmente por Roberts 
~\cite{roberts1959ewma} y consolidada como herramienta de control 
estadístico por Hunter~\cite{hunter1986ewma} resuelve este problema 
asignando pesos que decrecen geométricamente con la antigüedad de 
cada observación. Lucas y Saccucci~\cite{lucas1990ewma} ampliaron 
posteriormente este esquema con el análisis de sus propiedades de 
detección y con extensiones para señales con varianza variable. El 
baseline se actualiza en cada ciclo mediante la recurrencia
\begin{equation}
\mu_t = \alpha \cdot x_t + (1-\alpha) \cdot \mu_{t-1},
\label{eq:ewma}
\end{equation}
donde $\mu_t$ es el baseline estimado en el ciclo $t$, $x_t$ la 
observación actual, y $\alpha \in (0,1)$ el factor de suavizado que 
determina la velocidad de adaptación del estimador. Un valor de 
$\alpha$ cercano a uno produce un baseline altamente reactivo al 
último dato, mientras que un valor cercano a cero produce un 
baseline estable pero lento para adaptarse a cambios reales. El 
peso que una observación retiene tras $k$ ciclos es 
$\alpha \cdot (1-\alpha)^k$, lo que define una memoria efectiva 
finita a partir de la cual los datos antiguos se consideran 
olvidados~\cite{lucas1990ewma}.

La varianza se estima con el mismo principio exponencial 
(Exponentially Weighted Moving Variance, EWMV) mediante
\begin{equation}
{\sigma^2_t = (1-\alpha) \cdot \left[\sigma^2_{t-1} + \alpha \cdot (x_t - \mu_{t-1})^2\right],}
\label{eq:ewmv}
\end{equation}
donde $(x_t - \mu_{t-1})^2$ cuantifica la discrepancia entre la 
observación actual y el baseline anterior. Con el baseline y la 
varianza adaptativos, el Z-Score se calcula como
\begin{equation}
Z_t = \frac{x_t - \mu_{t-1}}{\sqrt{\sigma^2_{t-1}}},
\label{eq:zscore}
\end{equation}
contra el baseline y varianza del ciclo anterior, antes de 
actualizarlos con la observación actual. Este orden preserva la 
sensibilidad del estimador ante anomalías: si los estadísticos se 
actualizaran primero, la anomalía se absorbería parcialmente en el 
baseline antes de ser evaluada~\cite{montgomery2013spc}.

Bajo el supuesto de distribución aproximadamente normal en 
condiciones de operación estable, el Z-Score admite una 
interpretación probabilística directa. La probabilidad de que una 
observación produzca $|Z|>2$ bajo ese supuesto es 4.56\%, y la de 
que supere $|Z|>3$ es 0.27\%, lo que permite asignar una medida de 
rareza estadística a cada evaluación y establecer umbrales de 
severidad con significado probabilístico explícito 
~\cite{montgomery2013spc}.

\subsection{Arquitecturas cliente-servidor y patrones de integración}

La exposición de los datos procesados hacia la aplicación móvil se 
apoya en el estilo arquitectónico REST (Representational State 
Transfer) formulado por Fielding~\cite{fielding2000rest}. REST 
establece un conjunto de restricciones para sistemas distribuidos 
basados en HTTP, entre las cuales destacan la separación entre 
cliente y servidor, la ausencia de estado de sesión en el servidor, 
la identificación de recursos mediante URIs y la manipulación de 
esos recursos a través de los métodos estándar del protocolo HTTP 
~\cite{rfc9110}. La restricción \textit{stateless} implica que cada 
petición del cliente contiene toda la información necesaria para 
ser procesada, lo cual simplifica la escalabilidad horizontal del 
servidor pero requiere un mecanismo alternativo para mantener la 
autenticación del usuario entre peticiones, función que típicamente 
cumplen los tokens JWT~\cite{rfc7519} adjuntados en cada petición.

El modelo de petición y respuesta de HTTP es eficaz para consultas 
puntuales, pero presenta limitaciones en escenarios de monitoreo 
continuo donde el estado del sistema evoluciona de forma constante 
y el cliente necesita reflejar esos cambios sin retrasos 
significativos. El polling periódico, en el que el cliente repite 
una petición cada cierto intervalo, introduce sobrecarga de 
encabezados HTTP redundantes, consumo innecesario de ancho de banda 
y latencia variable entre la generación de un dato en el servidor y 
su visualización en el cliente~\cite{rfc9110}. El protocolo 
WebSocket~\cite{rfc6455} atiende este escenario mediante un canal 
bidireccional persistente establecido sobre una única conexión TCP 
tras un handshake inicial de actualización desde HTTP;@una vez 
abierto, tanto servidor como cliente pueden enviar mensajes en 
cualquier momento sin la sobrecarga de encabezados de HTTP, lo que 
habilita la notificación inmediata de eventos como alertas o 
actualizaciones de estado.

En sistemas compuestos por múltiples procesos que deben compartir 
información, el patrón \textit{Shared Database} descrito por Hohpe 
y Woolf~\cite{hohpe2003eip} propone una base de datos común como 
punto de acoplamiento entre procesos independientes. Cada proceso 
opera sobre su propio ciclo de vida, con sus propios recursos y 
perfil de carga, y accede a la base de datos para leer los datos 
que otros procesos produjeron o para escribir los suyos. Este 
patrón resulta apropiado cuando los procesos tienen perfiles de 
ejecución heterogéneos, ya que el aislamiento a nivel de proceso 
del sistema operativo garantiza que un fallo en uno de ellos no 
afecta la operación de los demás~\cite{silberschatz2018os}, a 
diferencia de una arquitectura monolítica en la que todas las 
responsabilidades comparten el mismo espacio de ejecución.

La organización interna de cada proceso se apoya en patrones de 
diseño clásicos documentados por Gamma et al.~\cite{gof1994}. Los 
patrones relevantes para este trabajo incluyen Strategy, que 
desacopla un algoritmo de su contexto de uso mediante una interfaz 
común; Mediator, que centraliza la comunicación entre componentes 
en un orquestador que conoce a todos los participantes y evita el 
acoplamiento directo entre ellos; y la inversión de control 
implementada mediante callbacks, que permite a un módulo invocar 
comportamiento definido externamente sin conocer los detalles de 
implementación~\cite{gof1994}. A estos se suman patrones 
específicos de sistemas concurrentes documentados por Schmidt et 
al.~\cite{schmidt2000posa2}, como Thread-per-Client, en el cual 
cada conexión entrante se atiende en un hilo dedicado para aislar 
los fallos individuales; y el patrón Productor-Consumidor descrito 
por Buschmann et al.~\cite{buschmann1996posa}, que coordina la 
transferencia de datos entre dos hilos con velocidades distintas 
mediante un buffer compartido protegido por exclusión mutua. La 
aplicación concreta de estos patrones se describe en la Sección 
~\ref{sec:desarrollo}.